\documentclass[12pt,a4paper]{article}

\usepackage[romanian]{babel}
\usepackage[T1]{fontenc}
\usepackage[utf8]{inputenc}
\usepackage{lmodern}
\usepackage{enumitem}

\title{Diagrama Conceptuală \\ Modelarea Conceptuală a Domeniului}
\author{}
\date{}

\begin{document}

\maketitle

\section{Modelarea Conceptuală}
Modelarea conceptuală reprezintă descrierea completă a domeniului (universului de
discurs), incluzând structuri, reguli, procese, entități, relații și constrângeri.
Rezultatul este \textbf{diagrama conceptuală}, care precede modelul logic (ERD).

\section{Componentele unei Diagrame Conceptuale}

\subsection*{1. Entități}
Obiecte sau concepte relevante pentru domeniu.

\textbf{Tipuri:}
\begin{itemize}[noitemsep]
    \item entități strong (independente);
    \item entități weak (dependente);
    \item entități asociative (pentru relații M:N);
    \item entități abstracte (în generalizare/specializare).
\end{itemize}

\subsection*{2. Atribute}
Proprietăți ale entităților.

\textbf{Tipuri:}
\begin{itemize}[noitemsep]
    \item simple, compuse, multivaloare, derivate;
    \item atribute-cheie (candidate, primare, alternative);
    \item domenii conceptuale (tipuri de date la nivel conceptual).
\end{itemize}

\subsection*{3. Relații}
Asocieri între entități.

\textbf{Tipuri:}
\begin{itemize}[noitemsep]
    \item unare (recursive);
    \item binare;
    \item ternare și n-are;
    \item relații de agregare și compoziție.
\end{itemize}

\subsection*{4. Cardinalități și Participare}
\begin{itemize}[noitemsep]
    \item min-max: (0,1), (1,1), (0,N), (1,N);
    \item participare totală sau parțială;
    \item constrângeri de rol (cine participă și cum).
\end{itemize}

\subsection*{5. Reguli de Business}
Reguli care guvernează domeniul, independente de implementare.

\textbf{Exemple:}
\begin{itemize}[noitemsep]
    \item un student nu poate fi înscris la mai mult de 6 cursuri;
    \item un profesor poate preda doar cursuri din specializarea sa.
\end{itemize}

\subsection*{6. Procese și Evenimente}
Descriu dinamica domeniului.

\textbf{Exemple:}
\begin{itemize}[noitemsep]
    \item proces: înscrierea unui student la un curs;
    \item eveniment: actualizarea notei;
    \item fluxuri: aprobarea unei cereri.
\end{itemize}

\subsection*{7. Generalizare / Specializare (ISA)}
\begin{itemize}[noitemsep]
    \item disjoint / overlapping;
    \item total / partial;
    \item moștenirea atributelor și relațiilor.
\end{itemize}

\subsection*{8. Constrângeri Conceptuale}
\begin{itemize}[noitemsep]
    \item integritate semantică;
    \item dependențe conceptuale;
    \item reguli de consistență între entități.
\end{itemize}

\subsection*{9. Perspective și Roluri}
\begin{itemize}[noitemsep]
    \item actori (student, profesor, administrator);
    \item responsabilități și interacțiuni;
    \item puncte de vedere asupra datelor.
\end{itemize}

\section{Mini-Diagramă Conceptuală (Textuală)}

\subsection*{Entități}
\begin{itemize}[noitemsep]
    \item Student(StudID, Nume, Email, DataNasterii)
    \item Curs(CursID, Denumire, Credite)
    \item Profesor(ProfID, Nume, Titlu)
    \item Inscriere(StudID, CursID, Data, Nota)
    \item TelefonStudent(StudID, NrTelefon)
\end{itemize}

\subsection*{Relații}
\begin{itemize}[noitemsep]
    \item Student (1,N) TelefonStudent
    \item Student (0,N) Curs prin Inscriere
    \item Profesor (1,N) Curs
    \item Angajat (1,N) seSubordonează Angajat (unară)
\end{itemize}

\subsection*{Reguli de Business}
\begin{itemize}[noitemsep]
    \item un student poate avea mai multe numere de telefon;
    \item un curs poate avea un singur profesor titular;
    \item un student poate fi înscris la maximum 6 cursuri simultan.
\end{itemize}

\subsection*{Generalizare}
\begin{itemize}[noitemsep]
    \item Angajat → Profesor, Administrator (disjoint, partial)
\end{itemize}

\subsection*{Procese}
\begin{itemize}[noitemsep]
    \item înscrierea la curs;
    \item evaluarea studentului;
    \item actualizarea datelor personale.
\end{itemize}

\end{document}
